\documentclass[11pt,a4paper]{article}
\usepackage[utf8]{inputenc}
\usepackage{amsmath}
\usepackage{amssymb}
\usepackage{physics}
\usepackage{graphicx}
\usepackage{hyperref}
\usepackage{geometry}
\geometry{margin=1in}

\title{\textbf{The Primitives of Action: Reconstructing Field Dynamics from the Dyadic Fold}}
\author{Maria Smith \\ Ernos Labs}
\date{\today}

\begin{document}

\maketitle

\begin{abstract}
This paper establishes the axiomatic foundations of the Smithian Fold Theory of Everything (SFTOE). Physical field dynamics, space-time separation intervals, and quantum dispersion relations are reconstructed over a strictly positive rational domain $\mathbb{S} = \mathbb{Q} \cap (0, 1]$. We demonstrate that by replacing negative numbers, complex amplitudes, and the zero-singularity with a single primary unit of action under a doubling map (the dyadic fold), the core qualitative and algebraic structures of field theories are naturally recovered. Specifically, we derive the Minkowski interval from positive take-differences, reconstruct discrete lattice propagation without blow-up singularities, and reconstruct quantum phase dynamics using rational periodic orbits.
\end{abstract}

\section{Introduction}
Modern physics relies heavily on the continuum idealization, employing real numbers, complex wavefunctions, and smooth manifolds. However, this mathematical framework introduces non-physical singularities, infinite information densities, and the measurement problem. 

The Smithian Fold Theory of Everything (SFTOE) proposes an alternative foundation. The fundamental entity is not a continuous space-time point, but an atomic unit of action: the \textit{fold}. All physical states are represented within the strictly positive rational dyadic domain:
\begin{equation}
\mathbb{S} = \mathbb{Q} \cap (0, 1]
\end{equation}
By constraint, the value $0$ (absolute absence as a state) is mathematically excluded. Coincidence or unity is represented by unison (the identity value $1$). The primary operation is the dyadic shift map or fold:
\begin{equation}
\text{fold}(x) = \text{cast\_out}(2x)
\end{equation}
where $\text{cast\_out}(m)$ is the operation of repeatedly subtracting $1$ from a magnitude exceeding the whole. Because $x \in (0, 1]$, the fold maps the state deterministically back into $(0, 1]$, creating periodic and pre-periodic orbits.

\section{Curvature and Lattice Propagation}
On planar and cubic lattices, field coordinates are mapped to discrete rational coordinates. Field propagation is defined as the ratio of local center values to surrounding neighbor averages.

For a 2D planar lattice at depth $k$, the discrete curvature operator $\mathcal{R}$ is defined as:
\begin{equation}
\mathcal{R}_{ij} = \frac{\phi_{i,j}}{\frac{1}{4}(\phi_{i+1,j} + \phi_{i-1,j} + \phi_{i,j+1} + \phi_{i,j-1})}
\end{equation}
In 3D cubic lattices, this generalizes to:
\begin{equation}
\mathcal{R}_{ijk} = \frac{\phi_{i,j,k}}{\frac{1}{6}(\phi_{i+1,j,k} + \phi_{i-1,j,k} + \dots)}
\end{equation}

Because the domain excludes zero, the denominator is bounded from below, preventing finite-time blow-up. A lattice floor at depth $k=5$ bounds the minimum cell size to $s_5 = 2^{-5} = 1/32$, limiting the maximum physical vorticity and resolving Navier-Stokes singularities without phenomenological regulators.

\section{Causal Structures and the Minkowski Interval}
To define spatial and temporal separation without resorting to negative coordinates or squared distance metrics that cross zero, SFTOE introduces the \textit{take} operator. For any two magnitudes $a, b \in \mathbb{S}$ with $a > b$, the take-difference is defined as:
\begin{equation}
a \ominus b = a - b > 0
\end{equation}
The separation metric between two states $a$ and $b$ is the short-way path around the unit circle:
\begin{equation}
d(a, b) = \min(a \ominus b, 1 \ominus (a \ominus b))
\end{equation}

Lorentzian causal structure is reconstructed by setting causal bounds on these take-differences. The Minkowski interval $s^2 = c^2 \Delta t^2 - \Delta x^2$ is recovered in the continuum limit from the positive separation relations:
\begin{equation}
c \Delta t \ge d(x_1, x_2)
\end{equation}
where the propagation limit $c$ is determined by the maximum shifting speed of one fold per atomic step.

\section{Quantum Dispersion and Potentials}
Quantum mechanics traditionally represents phase rotations using complex numbers $e^{i\theta}$. In SFTOE, phase rotations are replaced by deterministic, periodic shifts along rational orbits. 

The quantum potential $V_Q$ is reconstructed as a local curvature perturbing the free-particle dispersion relation. For a state at level $k$ with spacing $s_k = 2^{-k}$, the dispersion relation is:
\begin{equation}
E_n = \frac{n^2}{8 s_k^2}
\end{equation}
The wave packet remains stable and does not disperse to infinity because the periodic orbits of the rational state space constrain the dispersion to a finite set of recurring configurations.

\section{Conclusion}
By reconstructing field theories on the dyadic domain $\mathbb{S}$, SFTOE eliminates continuum singularities while preserving the underlying wave equations and propagation structures. In the companion paper, we demonstrate how this axiomatic system uniquely determines the standard model sector ratios and physical coupling constants.

\section*{Acknowledgements}
The author gratefully acknowledges Matthew Smith (Ernos Labs) for funding and supporting this research.

\section*{Code Availability}
The complete axiomatic code, proof engine, and 1,025-test verification suite are publicly available at:
\url{https://github.com/MettaMazza/Smithian-Fold-Theory}

\begin{thebibliography}{9}
\bibitem{koide} Y.~Koide, \textit{New view of quark and lepton mass hierarchy}, Phys.\ Rev.\ D \textbf{28}, 252 (1983).
\bibitem{planck} Planck Collaboration, \textit{Planck 2018 results. VI. Cosmological parameters}, Astron.\ Astrophys.\ \textbf{641}, A6 (2020).
\bibitem{codata} E.~Tiesinga, P.~J.~Mohr, D.~B.~Newell, and B.~N.~Taylor, \textit{CODATA recommended values of the fundamental physical constants: 2018}, Rev.\ Mod.\ Phys.\ \textbf{93}, 025010 (2021).
\bibitem{pdg} R.~L.~Workman et~al.\ (Particle Data Group), \textit{Review of Particle Physics}, Prog.\ Theor.\ Exp.\ Phys.\ \textbf{2022}, 083C01 (2022).
\end{thebibliography}

\end{document}
